% Part: applied-modal-logics
% Chapter: temporal-logic
% Section: introduction

\documentclass[../../../include/open-logic-section]{subfiles}

\begin{document}

\olfileid{aml}{tl}{int}

\olsection{ভূমিকা}

কালগত যুক্তিবিদ্যা এমন দাবি নিয়ে আলোচনা করে, যা ভবিষ্যতে সত্য হবে অথবা অতীতে সত্য ছিল। কালগত যুক্তিবিদ্যার প্রবর্তক হিসেবে আর্থার প্রায়রকে স্বীকৃতি দেওয়া হয়; তিনি একে \emph{কালরূপের যুক্তিবিদ্যা} বলেছিলেন। এখানে আমাদের আলোচনা প্রধানত প্রায়রের আদি মোডাল পদ্ধতি অনুসরণ করবে: বচনমূলক যুক্তিবিদ্যার মৌলিক কাঠামোয় কালগত অপারেটর যোগ করা। সেই কাঠামোতে দাবিগুলিকে সাধারণত কালরূপহীন বলে ধরা হয়।

উদাহরণ হিসেবে, বচনমূলক যুক্তিবিদ্যায় আমি বিজি নামের একটি কুকুরের কথা বলতে পারি, যে কখনও বসে, কখনও বসে না—কুকুরদের যেমন স্বভাব। বিজি বসে আছে এবং বিজি বসে নেই, এই দুই দাবি একসঙ্গে করা ধ্রুপদি যুক্তিবিদ্যায় স্ববিরোধী। কিন্তু দুটিই যে সত্য হতে পারে, তা স্পষ্ট; শুধু একই সময়ে নয়। ভাষায় কালগত অপারেটর যোগ করলে সেই দাবিটি বেশ সহজে প্রকাশ করা যায়। আরও, ``বিজি একটি খাবার বা একটি বল পাবে'' থেকে ``বিজি একটি খাবার পাবে অথবা বিজি একটি বল পাবে''—এ ধরনের অনুমানের বৈধতাও কালগত অপারেটর যোগ করে ব্যাখ্যা করা যায়।

তবে কালগত যুক্তিবিদ্যায় অনেক দার্শনিক প্রশ্ন ওঠে, যেগুলির জন্য আমরা একটি কাঠামোর বদলে অন্য কাঠামো বেছে নিতে পারি। যেমন, ভবিষ্যৎ-সম্পর্কিত আপতিক বচন এমন একটি ভবিষ্যৎ-বিষয়ক বক্তব্য, যা আবশ্যিকও নয়, অসম্ভবও নয়। ``রিচার্ড আগামীকাল মুদির দোকানে যাবে'' বললে আমরা এমন এক ঘটনা সম্পর্কে দাবি করি, যা এখনও ঘটেনি এবং যার সত্যমান নিয়ে বিতর্ক হতে পারে। বস্তুত, প্রথমেই সেই দাবিকে কোনো সত্যমান \emph{দেওয়া} যায় কি না, তা নিয়েও বিতর্ক আছে। আমরা যদি কঠোর নিয়তিবাদী হই, তবে সংশ্লিষ্ট ঘটনা ঘটার কথা যে সময়ে, তার আগেই বাক্যটি বাস্তবে সত্য অথবা মিথ্যা—এ ধারণা মেনে নিতে পারি; হয়তো আমরা তখনও তার সত্যমান জানি না। অন্যদিকে, আমরা ভবিষ্যৎকে সত্যিই উন্মুক্ত মনে করতে পারি, যেখানে ভবিষ্যৎ-সম্পর্কিত আপতিক বচনগুলির সত্যমান অনির্ধারিত।

দেখা যায়, সময়ের গঠন ও প্রকৃতি সম্পর্কে এমন অনেক অবস্থান কালগত যুক্তিবিদ্যার মডেল এবং কাঠামো বাছাইয়ের মধ্যেই নিহিত থাকে। যেমন, আমাদের প্রশ্ন হতে পারে, সময় রৈখিক, শাখাবিশিষ্ট, এমনকি চক্রাকার—কোন ধরনের মডেল তৈরি করা উচিত। আমাদের আরও স্থির করতে হতে পারে, কালগত মডেলের আদি ও অন্তবিন্দু থাকবে কি না এবং সময়কে বিচ্ছিন্ন মুহূর্ত দিয়ে, না একটি অবিচ্ছিন্ন পরিসর হিসেবে উপস্থাপন করা হবে।

\end{document}
